
\subsection{不同频率相互垂直简谐运动的合成：李萨如图形}

当两个相互垂直的简谐运动角频率不相等（$\omega_1 \neq \omega_2$）时，运动方程改写为：$$x = A_1 \cos(\omega_1 t + \phi_1),~~y = A_2 \cos(\omega_2 t + \phi_2)$$
只有当两个分振动的角频率（或频率、周期）之比严格满足整数比（有理数比）时，质点的总运动轨迹才会随时间演入重合，从而形成一条闭合且固定的几何曲线：$$\frac{\omega_1}{\omega_2} = \frac{\nu_1}{\nu_2} = \frac{m}{n} \quad (m, n \in \mathbb{Z}^+)$$在这种同频率成整数比的约束下，质点在平面内所描绘出的封闭稳定的运动轨迹曲线称为李萨如图形。
